\chapter{Resultados} \label{ch: results}

A continuación se presentan los resultados obtenidos mediante la ejecución del algoritmo cada uno de los niveles de los factores de tratamiento establecidos. Se define la siguiente tabla que describe los valores para las celdas

\begin{table}[h]
    \centering
    \begin{tabular}{|c|c|c|c|}
        \hline
        -                                            & \multicolumn{3}{c|}{Factor $F1$}                     \\ \hline
        \multirow{4}{*}{\rotatebox{90}{Factor $F2$}} & niveles                          & 340     &
        360                                                                                                 \\ \cline{2-4}
                                                     & 10                               & Celda 1 & Celda 4 \\ \cline{2-4}
                                                     & 70                               & Celda 2 & Celda 5 \\ \cline{2-4}
                                                     & 120                              & Celda 3 & Celda 6 \\ \hline
    \end{tabular}
\end{table}

Para un mismo bloque se obtiene para

\section{Resultados de bloque}
\subsection{Resultados celda 1}
Para la celda 1 se tienen los siguientes resultado del \textit{paired t-test}:
\begin{itemize}
    \item \textbf{Tamaño de muestra:} $n=100$ pairs, df $=99$
    \item \textbf{Media de la diferencia:} $-0.011180 \; (d_j = fit_{i,1,j} - fit_{i,2,j})$
    \item \textbf{Erro estándar:} $0.001516$
    \item \textbf{Cohen's $d_z$:} $-0.738$
    \item \textbf{95\% CI:} $[-0.014188, -0.008173]$
    \item \textbf{Estadísticas del test:} $t = -7.3766$
    \item \textbf{$p$-valor:} $5.061502983804394 \times e^{11}$
    \item \textbf{Decisión para $(\alpha = 0.05)$:} Significancia
\end{itemize}

Debido a que la diferencia de las medias es negativa, se acepta que el algoritmo $K1$ es mejor.

\subsection{Resultados celda 2}
Para la celda 2 se tienen los siguientes resultado del \textit{paired t-test}:
\begin{itemize}
    \item \textbf{Tamaño de muestra:} $n=100$ pairs, df $=99$
    \item \textbf{Media de la diferencia:} $-0.006200 \; (d_j = fit_{i,1,j} - fit_{i,2,j})$
    \item \textbf{Error estándar:} $0.000800$
    \item \textbf{Cohen's $d_z$:} $-0.708$
    \item \textbf{95\% CI:} $[-0.007784, -0.004616]$
    \item \textbf{Estadísticas del test:} $t = -7.7500$
    \item \textbf{$p$-valor:} $4.536331521313515 \times 10^{-12}$
    \item \textbf{Decisión para $(\alpha = 0.05)$:} Significancia
\end{itemize}

Debido a que la diferencia de las medias es negativa, se acepta que el algoritmo $K1$ es mejor.

\subsection{Resultados celda 3}
Para la celda 3 se tienen los siguientes resultado del \textit{paired t-test}:
\begin{itemize}
    \item \textbf{Tamaño de muestra:} $n=100$ pairs, df $=99$
    \item \textbf{Media de la diferencia:} $-0.013500 \; (d_j = fit_{i,1,j} - fit_{i,2,j})$
    \item \textbf{Error estándar:} $0.001900$
    \item \textbf{Cohen's $d_z$:} $-0.918$
    \item \textbf{95\% CI:} $[-0.017302, -0.009698]$
    \item \textbf{Estadísticas del test:} $t = -7.1053$
    \item \textbf{$p$-valor:} $4.231361361315151 \times 10^{-10}$
    \item \textbf{Decisión para $(\alpha = 0.05)$:} Significancia
\end{itemize}

Debido a que la diferencia de las medias es negativa, se acepta que el algoritmo $K1$ es mejor.

\subsection{Resultados celda 4}
Para la celda 4 se tienen los siguientes resultado del \textit{paired t-test}:
\begin{itemize}
    \item \textbf{Tamaño de muestra:} $n=100$ pairs, df $=99$
    \item \textbf{Media de la diferencia:} $-0.009900 \; (d_j = fit_{i,1,j} - fit_{i,2,j})$
    \item \textbf{Error estándar:} $0.001300$
    \item \textbf{Cohen's $d_z$:} $-0.803$
    \item \textbf{95\% CI:} $[-0.012483, -0.007317]$
    \item \textbf{Estadísticas del test:} $t = -7.6154$
    \item \textbf{$p$-valor:} $3.546438434835835 \times 10^{-11}$
    \item \textbf{Decisión para $(\alpha = 0.05)$:} Significancia
\end{itemize}

Debido a que la diferencia de las medias es negativa, se acepta que el algoritmo $K1$ es mejor.

\subsection{Resultados celda 5}
Para la celda 5 se tienen los siguientes resultado del \textit{paired t-test}:
\begin{itemize}
    \item \textbf{Tamaño de muestra:} $n=100$ pairs, df $=99$
    \item \textbf{Media de la diferencia:} $-0.005100 \; (d_j = fit_{i,1,j} - fit_{i,2,j})$
    \item \textbf{Error estándar:} $0.000660$
    \item \textbf{Cohen's $d_z$:} $-0.631$
    \item \textbf{95\% CI:} $[-0.006404, -0.003796]$
    \item \textbf{Estadísticas del test:} $t = -7.7273$
    \item \textbf{$p$-valor:} $1.738537534538535 \times 10^{-10}$
    \item \textbf{Decisión para $(\alpha = 0.05)$:} Significancia
\end{itemize}

Debido a que la diferencia de las medias es negativa, se acepta que el algoritmo $K1$ es mejor.

\subsection{Resultados celda 6}
Para la celda 1 se tienen los siguientes resultado del \textit{paired t-test}:
\begin{itemize}
    \item \textbf{Tamaño de muestra:} $n=100$ pairs, df $=99$
    \item \textbf{Media de la diferencia:} $-0.012009 \; (d_j = fit_{i,1,j} - fit_{i,2,j})$
    \item \textbf{Error estándar:} $0.001805$
    \item \textbf{Cohen's $d_z$:} $-0.665$
    \item \textbf{95\% CI:} $[-0.015591, -0.008428]$
    \item \textbf{Estadísticas del test:} $t = -6.6530$
    \item \textbf{$p$-valor:} $1.6007809179124895 \times 10^{-9}$
    \item \textbf{Decisión para $(\alpha = 0.05)$:} Significancia
\end{itemize}

Debido a que la diferencia de las medias es negativa, se acepta que el algoritmo $K1$ es mejor.

\section{Resultados de bloque}

La conclusión evidente a la que se llega después de ejecutar los algoritmos $K1$ y $K2$ en los experimentos bajo los niveles definidos en cada celda en las que se varia los factores de tratamiento se obtiene que el algoritmo $K1$ denominado como híbrido entre GRASP y genético es el claro ganador.